% \Fidel{TODO:
% \begin{itemize}
%     \item Sections 2-5 aren't very QEC-specific as we have a general framework. Maybe add a subsection (or just include in the section 6.0 introduction) discussing how SPPs can be used to model correlated noise in QEC specifically, and why we model the way we do in sections 6 and 7.
%     \item Justify that although our framework is general enough to capture correlations at the circuit level (we have freedom to choose where correlated errors occur throughout the circuit, since the twirl is applied at the gate level), we focus on phenomenological noise models, where errors occur at the granularity of QEC cycles, because (a) circuit-level Pauli noise can be mapped to correlated phenomenological noise, and (b) we are more interested in behaviour of correlations over multiple QEC cycles.
%     \item We included a decent amount of tensor network formalism in sections 2-5, which we don't end up explicitly using in sections 6 and 7. For our particular application for efficient large-scale circuit-level simulations, they are more intermediate steps to building efficient HMM models derived from the microscopic dynamics. We should address this.
% \end{itemize}
% }
% \begin{itemize}
%     \item Overview: Previous sections developed the general SPP framework. For the remainder of the paper, we focus on applications to QEC for modelling general spatiotemporal correlations in Pauli noise models occurring during QEC cycles.
%     \item Motivation (brief): Realistic noise in quantum devices exhibits non-trivial spatiotemporal correlations, which can significantly impact the performance of QEC codes and decoders. Therefore, it is crucial to have efficient and accurate models for simulating QEC under correlated noise. Maybe briefly mention previous work on the detrimental impact of correlated noise on QEC, and also experimental observations of correlated noise and their impact.
%     \item Posit SPPs as the general spatiotemporal correlated Pauli noise model for QEC. Can be used for correlated circuit-level noise models, since RC is applied at the gate level, but is also well-suited for phenomenological noise models at the QEC cycle level, which `live closer' to the decoder, and syndrome based characterisation methods. We have a choice on the `granularity' of the SPP model for QEC.
%     \item Mention we make use of the analytical tools developed throughout the paper.
%     \item Hint of other QEC related applications for SPPs, e.g., learning and tomography, tailoring decoding algorithms to correlated noise. Refer to discussion section for more details. Also immediate potential for other immediate extensions, such as different noise models, and different QEC codes and circuits (e.g., flouqet and QLDPC) all compatible with our theoretical and numerical framework.
%     \item Justify the bridge from correlated circuit-level Pauli noise to correlated phenomenological noise at the QEC cycle level beyond simplification. First, circuit-level Pauli noise can be mapped to correlated phenomenological Pauli. Second, we are more interested in the behaviour of correlations over multiple QEC cycles. Third, related to the previous, is that the emergent phenomenological Pauli correlations at the QEC cycle level of syndrome measurements, have major implications for decoding performance and characterisation. For example, ML based decoders like AlphaQubit 2 increase logical performance by exploiting correlations in phenomenological noise.
%     \item Related to the above, articulate and justify our transition from tensor network formalism to HMMs. Although the tensor network methods are universally applicable, for the purposes of studying large-scale QEC simulations under correlated noise, tensor networks are more of a stepping stone to efficient HMM models derived from the microscopic dynamics. This allows us to efficiently simulate large-scale QEC codes under non-trivial correlated noise models.
%     \item Summarise and explain our choices of models in sections 6 and 7. This is more a generative modelling paper than predictive modelling. The temporal storm model in this section is a simple, analytically tractable 1D model that captures essential features of temporal correlations in noise. The QCA model in the next section is a more complex 2D model that captures both spatial and temporal correlations, derived from the microscopic SE quantum dynamics, has interesting statistical mechanical properties though is only defined by a few parameters. Note that our example employs a mix of independent circuit-level noise and correlated inter-round errors modelled by the SPP.
%     \item Really try to sell both the SPP framework and the models. Though this will also be emphasised in the introduction. E.g., go back to the bridge story of microscopic SE dynamics on a device all the way to the effective Pauli noise model relevant for QEC. Cite previous works on complex noise models for QEC, which often have a mix of limitations such as: (a) ad hoc noise models that either don't derive from the microscopic dynamics of have strong assumptions like perfect syndrome extraction (b) limited correlations, e.g., only spatial (bath resets in between cycles) or short-range (Gaussian) correlations (c) limited to small codes in contrast to the hundreds of qubits needed for FTQC (d) limited to specific codes or decoders, for instance, previous studies rely on statistical mechanical mappings that only work for specific codes, noise models, and decoders (e) focus on threshold estimates, which have fallen out of favour compared to logical performance under threshold (teraquop footprint is king).
%     \item Maybe also try to sell the QCA model further as complex and interesting in its own right, beyond just being a correlated noise model for QEC.
%     \item I realise that this is basically becoming a second introduction, which is fine. The job of this section is to tie the general SPP framework to a specific QEC application (modelling correlated noise), but where the actual introduction should be more of a high level sales pitch, this should be more detailed and explanatory. Want this to section to be concise and punchy as possible and not include repetitive details.
% \end{itemize}

% Although the framework resolves correlations at the level of individual gates, for large-scale QEC simulations we adopt a hybrid description of independent gate-level noise combined with inter-round correlations modelled by an SPP at QEC cycle level. This choice is motivated by several considerations. First, any Pauli noise model on a (Clifford) syndrome extraction circuit can be coarse-grained to a correlated Pauli error model on data and ancilla qubits at the QEC cycle level. Second, the correlations that most strongly impact QEC performance often manifest over multiple rounds, rather than within a single cycle. Third, decoders and certain characterisation protocols operate on coarse-grained objects on the level of syndrome outcomes. As a result, phenomenological SPPs provide a natural interface between microscopic descriptions and practical QEC tools.

% Realistic noise in quantum devices exhibits nontrivial correlations across both space and time, which can substantially alter QEC performance relative to memoryless baselines. Many existing treatments therefore introduce a correlated noise model tailored to the phenomenon of interest, frequently relying on simplifying assumptions about microscopic structure, correlation range, or the syndrome extraction setting. Our goal here is to not catalogue models, but to emphasise that a practical QEC-facing description of noise should (i) capture correlations at a level relevant to syndrome processing and decoding, while (ii) remain compatible with general underlying (Markovian or non-Markovian) dynamics.

% SPPs provide a convenient language for this purpose. They encode correlated stochastic Pauli faults across space and time, and can be connected operationally and practically to general noise via the multi-time Pauli twirl implemented by randomised compiling. Formally, the multi-time twirl maps the process tensor description of a noisy stabiliser circuit to an SPP that assigns a joint distribution over Pauli faults on gates and idle intervals throughout the computation. This forms the concrete bridge from microscopic dynamics to a correlated Pauli description suitable for QEC analysis.
\clearpage
\section{Applications for quantum error correction}\label{sec:temporalstormmodel}

\subsection{Spatiotemporal Pauli processes as QEC noise models}
Up to this point, we have developed a general framework for spatiotemporal Pauli processes (SPPs). In principle, an SPP description of circuit noise can be obtained from a tomographically reconstructed process tensor~\cite{whiteDemonstrationNonMarkovianProcess2020,whiteUnifyingNonMarkovianCharacterization2025}. In practice, it is more useful to characterise an SPP directly at a chosen level of granularity---either at the circuit level, or phenomenologically at the level of QEC cycles. For the remainder of the paper, we focus on SPPs as generative noise models for QEC; broader uses, including decoding and noise mitigation, are discussed in \cref{sec:discussion}.

Correlations in realistic device noise can substantially alter QEC performance relative to memoryless baselines. Our goal here is to not catalogue and deconstruct existing models, but to emphasise that a practical QEC-facing description of noise should (i) capture correlations at a level relevant to syndrome processing and decoding, while (ii) remain compatible with general underlying (Markovian or non-Markovian) dynamics. SPPs provide a convenient language for this purpose. Under randomised compiling, the multi-time Pauli twirl maps noisy stabiliser circuit dynamics to an effective SPP---a joint distribution over Pauli fault across space and time.

Although the SPP framework resolves correlation at the level of individual operations, large-scale simulations and decoder interfaces often motivate a coarser representation. Accordingly, in our numerical studies we adopt a hybrid description: we treat within-round circuit noise using a standard independent gate-level model, while modelling inter-round memory effects by an SPP defined at the QEC cycle level. This choice reflects three pragmatic considerations. First Pauli noise on a Clifford syndrome extraction circuit can be coarse-grained to an effective correlated Pauli error model on data and ancilla qubits at the end or start of each cycle. Second, the correlations that most strongly impact long-time performance are frequently those that manifest across multiple rounds. Third, decoders and many characterisation protocols naturally consume coarse-grained syndrome data, making a cycle-level SPP a convenient interface between microscopic noise and practical QEC tools.

In this section and the next, we exploit this bridge from microscopic dynamics, through SPPs, to efficiently simulable HMMs in two concrete settings. In \cref{sec:temporalstormmodel}, we introduce a simple one-dimensional SPP with tunable spectral gap (equivalently, tunable correlation length) while keeping marginal error rates fixed, enabling a controlled study of how temporal correlations can impact surface code memory and stability experiments. In \cref{sec:qcamodel}, we turn to a genuinely spatiotemporal two-dimensional SPP derived from a quantum cellular automaton model of environmental dynamics, which generates rich correlation structure and statistical-mechanical behaviour while remaining parametrically simple.

\subsection{Temporal storm noise model}
\begin{figure}[t]
    \centering
    \begin{subfigure}[b]{0.49\textwidth}
      \centering
      \valignboxnopad{\begin{tikzpicture}
	\begin{pgfonlayer}{nodelayer}
		\node [style=smallscaledcircleorange] (0) at (0.5, 0) {\(\mathbf{A}_0\)};
		\node [style=smallscaledcircleorange] (1) at (1.625, 0) {\(\mathbf{A}_1\)};
		\node [style=smallscaledcircleorange] (2) at (3.05, 0) {\(\mathbf{A}_k\)};
		\node [style=none] (3) at (0.5, -0.5) {};
		\node [style=none] (4) at (1.625, -0.5) {};
		\node [style=none] (5) at (3.05, -0.5) {};
		\node [style=none] (6) at (2.375, 0) {};
		\node [style=none] (8) at (0.5, -0.625) {\({}_{x_0}\)};
		\node [style=none] (9) at (1.625, -0.625) {\({}_{x_1}\)};
		\node [style=none] (10) at (3.05, -0.625) {\({}_{x_k}\)};
		\node [style=none] (11) at (1.075, 0.2) {\({}_{\mu_1}\)};
		\node [style=none] (12) at (2.2, 0.2) {\({}_{\mu_2}\)};
		\node [style=none] (14) at (3.625, 0) {};
		\node [style=none] (15) at (4.375, 0) {};
		\node [style=smallcircleblue] (16) at (5.125, 0.625) {\(_{S=0}\)};
		\node [style=smallcirclered] (17) at (6.125, -0.625) {\(_{S=1}\)};
		\node [style=none] (18) at (6.25, 0.375) {\(_{\Gamma_{0\to 1}}\)};
		\node [style=none] (19) at (5.125, -0.375) {\(_{\Gamma_{1\to 0}}\)};
		\node [style=none] (20) at (5.375, 1.25) {\(_{\Gamma_{0\to 0}}\)};
		\node [style=none] (21) at (6, -1.25) {\(_{\Gamma_{1\to 1}}\)};
	\end{pgfonlayer}
	\begin{pgfonlayer}{edgelayer}
		\draw (5.center) to (2);
		\draw (4.center) to (1);
		\draw (3.center) to (0);
		\draw [style=fusedwire] (1) to (0);
		\draw [style=fusedwire] (6.center) to (1);
		\draw [style=doublearrow] (14.center) to (15.center);
		\draw [style=arrow, bend left] (16) to (17);
		\draw [style=arrow, bend left] (17) to (16);
		\draw [style=arrow, in=165, out=105, loop] (16) to ();
		\draw [style=arrow, in=-15, out=-75, loop] (17) to ();
		\draw [style=fuseddotted] (6.center) to (2);
	\end{pgfonlayer}
\end{tikzpicture}
}
      \caption{}
    \end{subfigure}
    \begin{subfigure}[b]{0.49\textwidth}
      \centering
      \valignboxnopad{\begin{tikzpicture}
	\begin{pgfonlayer}{nodelayer}
		\node [style=circboxwhite] (0) at (1.625, 0.5) {\(\mathcal{G}\)};
		\node [style=circboxblue] (1) at (2.3, 0.5) {\(I\)};
		\node [style=circboxwhite] (2) at (2.975, 0.5) {\(\mathcal{G}\)};
		\node [style=circboxred] (3) at (3.65, 0.5) {\(X\)};
		\node [style=circboxwhite] (4) at (4.325, 0.5) {\(\mathcal{G}\)};
		\node [style=circboxred] (5) at (5, 0.5) {\(Y\)};
		\node [style=circboxwhite] (6) at (5.675, 0.5) {\(\mathcal{G}\)};
		\node [style=circboxblue] (7) at (6.35, 0.5) {\(I\)};
		\node [style=smallcircleblack] (11) at (0.5, 0.525) {};
		\node [style=smallcirclewhite] (12) at (1, 1.025) {};
		\node [style=smallcirclewhite] (13) at (0, 1.025) {};
		\node [style=smallcirclewhite] (14) at (0, 0.025) {};
		\node [style=smallcirclewhite] (15) at (1, 0.025) {};
		\node [style=none] (16) at (0, 1.025) {};
		\node [style=none] (17) at (1, 1.025) {};
		\node [style=none] (18) at (1, 0.025) {};
		\node [style=none] (19) at (0, 0.025) {};
		\node [style=none] (20) at (1.225, 0.5) {};
		\node [style=smallercircleblue] (22) at (2.3, 1.5) {0};
		\node [style=smallercirclered] (23) at (3.65, 1.5) {1};
		\node [style=smallercirclered] (24) at (5, 1.5) {1};
		\node [style=smallercircleblue] (25) at (6.35, 1.5) {0};
	\end{pgfonlayer}
	\begin{pgfonlayer}{edgelayer}
		\draw [style=fillblue] (19.center)
			 to (16.center)
			 to (17.center)
			 to (18.center)
			 to cycle;
		\draw [style=fuseddotted] (11) to (20.center);
		\draw [style=fusedwire] (20.center) to (0);
		\draw [style=arrow, bend left] (22) to (23);
		\draw [style=arrow, bend left] (23) to (24);
		\draw [style=arrow, bend left] (24) to (25);
		\draw [style=arrowthick] (22) to (1);
		\draw [style=arrowthick] (23) to (3);
		\draw [style=arrowthick] (24) to (5);
		\draw [style=arrowthick] (25) to (7);
		\draw [style=fusedwire] (0) to (1);
		\draw [style=fusedwire] (1) to (2);
		\draw [style=fusedwire] (2) to (3);
		\draw [style=fusedwire] (3) to (4);
		\draw [style=fusedwire] (4) to (5);
		\draw [style=fusedwire] (5) to (6);
		\draw [style=fusedwire] (6) to (7);
	\end{pgfonlayer}
\end{tikzpicture}
}
      \caption{}
    \end{subfigure}
    \caption{Temporal storm noise model. A two-state latent environment \(s_t \in \{0,1\}\) (calm, storm) evolves as a Markov chain and modulates the Pauli label \(x_t \in \{I, X, Y, Z\}\) emitted at each time step, yielding temporally correlated Pauli noise with fixed single-round marginals. (a) Equivalent representations of the same one-dimensional SPP as an MPS and as a two-state HMM with emission probabilities \(q_{s_t}^{(x)}\). (b) Cycle-level noise schematic used for our QEC simulations. Each qubit (data and ancilla) couples to an independent two-state environment; \(\mathcal{G}\) denotes a chosen granularity of circuit operations (from a single gate to an entire syndrome extraction subcircuit).}
    \label{fig:temporal-storm-overview}
\end{figure}

We now instantiate the general SPP formalism in a one-dimensional, analytically tunable setting, which we refer to as the \emph{temporal storm model}. Each system qubit is coupled to its own two-level environment with classical basis states $\{\ket{0},\ket{1}\}$, representing \emph{calm} and \emph{storm} conditions, respectively.

The environment undergoes a time-homogeneous classical Markov chain with transition probabilities $\Gamma_{s\to s'}$ between calm and storm (\cref{fig:temporal-storm-overview}). The system-environment interaction is modelled as a quantum channel conditioned on the environment state. When the environment \(E\) is in state $\ket{s'}$, the system \(S\) evolves according to
\begin{equation}
    \mathcal{E}_{s'}[\rho_S]
    = \sum_x q^{(x)}_{s'} \,\sigma^{(x)} \rho_S \sigma^{(x)},
\end{equation}
where $\sigma^{(x)}$ denotes the Pauli operators indexed by $x$, and $q^{(x)}_{s'}$ are probabilities satisfying \(\sum_x q^{(x)}_{s'} = 1\). The joint system-environment evolution over a single time step is thus described by the CPTP map
\begin{equation}
    \Phi[\rho_{SE}]
    = \sum_{s,s',x} K^{(x)}_{s's} \,\rho_{SE}\, K^{(x)\dagger}_{s's},
    \qquad
    K^{(x)}_{s's}
    = \sqrt{\Gamma_{s\to s'}\,q^{(x)}_{s'}}\;
      \ketbra{s'}{s}_E \otimes \sigma^{(x)}_S.
\end{equation}
Because the environment dynamics are classical and the conditional system channels \(\mathcal{E}_{s'}\) are Pauli, the corresponding process tensor is invariant under the multi-time Pauli twirl in \cref{eq:twirl-upsilon}; in other words, this construction directly specifies an SPP, and may be viewed as the twirled image of a family of compatible underlying process tensors.

Labelling the environment transition probabilities as $a = \Gamma_{0\to 1}$ and $b = \Gamma_{1\to 0}$, the corresponding transfer operator $T$ and SPP matrices $A^{(x)}$ are
\begin{equation}
    T
    =
    \begin{pmatrix}
        1 - a & a \\
        b & 1 - b
    \end{pmatrix},
    \qquad
    A^{(x)} = T\,\mathrm{diag}\big(q_0^{(x)}, q_1^{(x)}\big).
\end{equation}
This is precisely a two-state HMM where \(T\) describes the evolution of the latent state, and the Pauli emission label \(x_t\) is drawn from \(q_{s_t}^{(x)}\). The spectrum of $T$ is directly obtained as \(\lambda_1 = 1, \lambda_2 = 1 - a - b\). The subleading eigenvalue \(\lambda_2\) sets the spectral gap \(\Delta\) and correlation length \(\xi\),
\begin{equation}\label{eq:storm-spectral-gap}
    \Delta = a + b, \quad \xi = -\frac{1}{\ln(1-a-b)},
\end{equation}
where we assume \(0 \le a\), \(b \le 1\) with \(a + b < 1\), so that \(\lambda_2 \in (0, 1)\) and \(\xi > 0\).

The left and right eigenvectors associated with $\lambda_1$, corresponding to the left and right stationary states, are
\begin{equation}\label{eq:storm-stationary-states}
    \bra{l} =
    \begin{pmatrix}
        \pi_0 & \pi_1
    \end{pmatrix},
    \qquad
    \ket{r} =
    \begin{pmatrix}
        1 \\
        1
    \end{pmatrix},
\end{equation}
where \(\pi_0 = b/(a+b)\) and \(\pi_1 = a/(a+b)\) are the stationary calm and storm fractions, respectively. These immediately lead to the marginal Pauli error rates,
\begin{equation}\label{eq:storm-marginals}
    \bar{p}^{(x)}
    = \pi_0 q_0^{(x)} + \pi_1 q_1^{(x)} = \frac{bq_0^{(x)} + aq_1^{(x)}}{a+b}.
\end{equation}
By solving \cref{eq:storm-spectral-gap,eq:storm-marginals} for the transition rates \(a\) and \(b\), we construct a family of temporally correlated noise models with tunable spectral gap \(\Delta\) (or \(\xi\)) while maintaining fixed marginal error rates \(\bar{p}^{(x)}\). This decoupling of environmental memory from average noise strength underpins the parameter sweeps in the surface code simulations that follow.

\subsection{Methodology}

\begin{itemize}
    \item Briefly summarise surface code stability vs surface code memory, and why it's important for temporally correlated noise.
\end{itemize}

\subsection{Results}
\begin{figure}
    \centering
    \includegraphics[width=1\textwidth]{figures/experiment1.pdf}
    \caption{}
\end{figure}

\begin{figure}
    \centering
    \includegraphics[width=1\textwidth]{figures/experiment2.pdf}
    \caption{}
\end{figure}

